\section{Data Generation, Quality, Lineage and Privacy}\label{sec:data}

\subsection{Public sources and the generator}
District boundaries, codes and the 2022 census population by age band and sex come from the HDX COD-AB and
COD-PS datasets~\cite{hdx2025codab}; daily precipitation, temperature and humidity at each district centroid
for 2021--2025 come from the Open-Meteo archive~\cite{zippenfeld2023openmeteo}. Open-Meteo allows about 5,000
weighted calls an hour and 64 districts over five years cost about 8,400, so the download is throttled and the
processed reference files (0.6\,MB) are committed. A clean clone therefore skips a two-hour download.

The generator (\texttt{generator/}) writes the three synthetic sources: encounters (S4) as one CSV folder per
report date, a register of 736 facilities (S5) and daily stock snapshots (S6) as JSON. It is vectorised numpy,
seeded per (seed, stream, district, year), so the same arguments give byte-identical output whatever the number
of worker processes. Rates are per head of population, so the row count follows the census rather than a
constant: SF=1 gives about 6.2 million encounters a year, or about 31 million for 2021--2025.

Each disease has its own annual cycle. Dengue peaks in early October, because OpenDengue's national counts peak
in October in 2022 and 2024, later than ``follows the monsoon'' suggests. Dengue responds to rainfall and
temperature four to eight weeks earlier, diarrhoea spikes after extreme rainfall, malaria is concentrated in
the three Hill Tracts districts and Cox's Bazar, and every count has negative-binomial noise. Reporting delay is
a discretised negative binomial with mean 3 days truncated at 21, and 65\% of Friday and Saturday reports move
to Sunday (Bangladesh's weekend). The SF=0.05 run injected 35 outbreaks\prov{} of different shape and size,
including a slow-growing one and a nationwide novel respiratory surge, and two reporting artefacts: a facility
that goes silent for three weeks and then back-fills. They are written only to \texttt{ground\_truth/}, which
the detection code never reads (a test fails if it does).

The generator also plants the defects the quality rules must catch: 2\% repeated \texttt{encounter\_id},
0.5\% orphan \texttt{facility\_id}, 0.5\% invalid age, 0.5\% invalid sex, 0.3\% malformed ICD-10 codes and
0.3\% onsets after the report date. A 14-day incremental drop (\texttt{make drop}) uses the same layout and
carries records with onsets up to 21 days before the drop.

\subsection{Bronze: keep what arrived}
\texttt{pipeline/bronze/ingest.py} finds files by glob, never by name, date or count. It computes each file's
SHA-256 and skips any file whose checksum is already in the ingest log, so a re-run reads nothing and a new
drop loads only its own files. Every value of the contract fields is kept as a string, exactly as written, and
each row gets the source, its file, the checksum, the batch ID and the load time (ING-1). One run is one batch;
batch IDs start with a UTC timestamp, so they sort in load order.

\subsection{Silver: rules, quarantine and de-identification}
The rule catalogue (\texttt{pipeline/quality/rules.py}) has 14 rules for encounters and 7 each for facilities
and stock (DQ-1). Each rule is a Spark column expression that is true when a row breaks it, with a stable reason
code, a severity (quarantine or warn) and the requirement it serves. The six planted defects map one to one onto
six rules:
\begin{itemize}
\item \texttt{DQ\_DUP\_ENCOUNTER} and \texttt{DQ\_ORPHAN\_FACILITY};
\item \texttt{DQ\_INVALID\_AGE} and \texttt{DQ\_INVALID\_SEX};
\item \texttt{DQ\_ONSET\_AFTER\_REPORT} and \texttt{DQ\_BAD\_ICD10}.
\end{itemize}
The rest cover nulls, types, long delays, future dates and outcome values. Spark~4 runs in ANSI mode, where a plain cast of a bad
string throws, so every rule reads values through \texttt{try\_cast} and an unparseable value is a defect of
its own. A failing row goes to \texttt{silver.quarantine} once, with the list of every code it broke. In the
first SF=0.05 load, 66,979 of 1,681,896 encounter rows (4.0\%)\prov{} were quarantined. Each load writes one
row of rule pass rates to \texttt{gold.dq\_report}, which the dashboard shows (DQ-2).

Before a clean row is stored, \texttt{patient\_id} becomes an HMAC-SHA256 under a key held only in
\texttt{.env} (PRV-1). Spark has no HMAC function and a Python UDF over 25 million rows is slow, so the HMAC is
built from Spark's native \texttt{sha2} with the RFC~2104 inner and outer pads. A plain salted hash would not be
enough: national IDs are few enough to enumerate, so the secret key is the only protection. The code refuses to
run with a missing key or the placeholder from \texttt{.env.example}. Exact age becomes an age band.

\subsection{Idempotent incremental loads (ING-4)}
Clean rows are merged into \texttt{silver.encounters} with a Delta MERGE on \texttt{encounter\_id}. A matched
row is replaced only when the incoming \texttt{report\_date} is later. That one condition gives both required
properties: re-running a batch matches every row with an equal date and changes nothing, and a resubmitted or
late record replaces the old copy and moves to its own onset week. The MERGE is limited to the epi-years in the
batch (plus one either side), so partition pruning keeps it from scanning the whole table.

Silver records the weeks each batch changed: the onset weeks of new and replaced rows and the old weeks of
replaced rows. Gold rebuilds only those weeks, plus the week after each, with a Delta \texttt{replaceWhere}.

\subsection{Gold, lineage and patient de-duplication}
\texttt{gold.cases\_weekly} counts cases, admissions, deaths and lab confirmations per district, disease and
ISO week of \emph{onset}. Counting by onset rather than report date is what keeps the back-filled weeks of the
silent facility from looking like an outbreak. Every district-week with at least one reporting facility gets a
row for all four diseases, so a missing row means ``nothing reported'' rather than zero, which the detector
uses to treat silent weeks as missing (AN-3). Each row carries the sorted bronze batch IDs of the records it
counts, which link through the ingest log to file names and checksums (DQ-3).

Cross-facility patient de-duplication (DQ-4) flags a record as a repeat when the same patient key, disease,
sex and compatible age band appear at another facility within seven days. On SF=0.05 it flagged 441 pairs, of
which 440 were true repeats (precision 0.998)\prov{}. Recall within the window was only 0.11\prov{}, because the
generator draws age and sex afresh for a repeat visit, so most planted repeats fail the attribute check. We
kept the strict rule: a false merge would hide a real case.

\subsection{Hidden-drop drill}
\texttt{scripts/hidden\_drop\_drill.sh} rehearses the instructor's test: clone, \texttt{make setup},
\texttt{make all}, a new 14-day drop with an unseen seed, the incremental steps, then the same steps again. On
11 October at SF=0.05 it passed every check\prov{}: the drop was one new bronze batch, silver gained 8,292 rows,
the data moved from 31 December 2025 to 14 January 2026, the re-run changed nothing and left zero duplicate IDs
and zero duplicate gold keys, and the export was 9.64\,MB with no count from 1 to 4. It took 22 minutes,
including setup.
